Status as
of
2026-10-05.
This
document
is the
single
source for
the
paper's
Background
section
and the PI
presentation
on this
topic. It
describes,
step by
step, how
we tested
GCAM-KAIST's
native
non-CO2
emissions,
and why
the
project
builds
bottom-up
emission
factors
instead.

It builds
on the
audit in
{\texttt{gcam\_\allowbreak{}kaist\_\allowbreak{}nonco2\_\allowbreak{}audit.\allowbreak{}md}}
(cited
below as
``audit
§\ldots{}''),
which
documents
how
GCAM-KAIST
stores and
reports
non-CO2
emissions.
Facts
established
there are
cited, not
repeated.

\textbf{Framing.}
Almost
everything
described
here is
GCAM's
standard
non-CO2
approach,
which
GCAM-KAIST
inherits
from stock
GCAM v9
(audit
§B5.2).
The
critique
is of that
approach
as applied
to
policy-driven
technology
change in
Korea, not
of choices
made by
the KAIST
team.

\textbf{Plain-English
sections}
are marked
\plainmark{}.
Technical
sections
follow
each one.

\hypertarget{summary}{%
\subsection*{0. Summary \plainmark{}}\addcontentsline{toc}{subsection}{0. Summary}\label{summary}}

The
question
was
whether we
can take
GCAM-KAIST's
own
non-CO2
emissions
(NOx, SO2,
VOCs, NH3,
CO, BC)
and use
them to
estimate
the
air-quality
and health
effects of
Korea's
net-zero
pathway.

We can't.
The core
reasons:

\begin{enumerate}
\def\labelenumi{\arabic{enumi}.}
\tightlist
\item
  \textbf{GCAM-KAIST's
  close
  agreement
  with
  Korea's
  national
  inventory
  (CAPSS)
  is not
  evidence
  of
  accuracy.}
  The
  base-year
  numbers
  \emph{are}
  inventory
  numbers:
  CAPSS
  totals
  passed
  through
  CEDS
  into
  GCAM's
  calibration.
  Comparing
  them
  with
  CAPSS
  compares
  CAPSS
  with
  itself.
\item
  \textbf{The
  emission
  factors
  are a
  by-product
  of that
  calibration,
  not
  measurements.}
  GCAM
  divides
  each
  inventory
  total by
  its own
  modelled
  activity.
  Different
  technologies
  inside
  the same
  total
  therefore
  get the
  same
  factor,
  and the
  factors
  swing
  from one
  base
  year to
  the
  next.
\item
  \textbf{After
  2021 the
  factors
  don't
  respond
  to
  policy.}
  They
  stay
  frozen
  or
  follow a
  GDP-based
  curve
  that's
  the same
  in both
  scenarios.
  Technologies
  with no
  emissions
  tag
  count as
  zero,
  and the
  net-zero
  scenario
  shifts
  activity
  toward
  exactly
  those
  technologies.
\item
  \textbf{These
  limits
  change
  the
  policy
  answer.}
  Keeping
  GCAM-KAIST's
  activity
  and
  swapping
  only the
  emission
  factors,
  the
  projected
  net-zero
  change
  in steel
  SOx
  flips
  from
  −88\% to
  +19\%,
  and the
  power-sector
  NOx and
  SOx
  avoided is 2.6--5.2× larger.
\end{enumerate}

\textbf{What
we do
instead:}
keep
GCAM-KAIST's
\emph{activity},
and
replace
its
\emph{emission
factors}
with
bottom-up,
technology-specific
Korean
factors,
validated
against
CAPSS in
the base
years.

\hypertarget{the-question}{%
\subsection*{1. The question \plainmark{}}\addcontentsline{toc}{subsection}{1. The question}\label{the-question}}

The NZK
health
analysis
needs
emissions
that
respond to
what a
decarbonization
pathway
actually
changes:
which
fuels are
burned, in
which
technologies,
with what
pollution
controls.
Those
emissions
then go
into InMAP
and a
health
model.

GCAM-KAIST
is the
project's
source of
projected
activity:
how much
electricity
each power
technology
generates,
how much
steel each
production
route
makes, how
many
kilometres
each
vehicle
type
drives,
and so on.
It also
reports
its own
non-CO2
emissions.
The
obvious
shortcut
would be
to use
those
directly.

Our first
comparison
against
CAPSS
looked
encouraging:
some
categories
matched
closely.
So the
question
became:
\textbf{does
that
agreement
mean
GCAM-KAIST's
native
non-CO2
emissions
can be
trusted
for policy
analysis?}

\hypertarget{how-gcam-represents-non-co2-emissions}{%
\subsection*{2. How GCAM represents non-CO2 emissions \plainmark{}}\addcontentsline{toc}{subsection}{2. How GCAM represents non-CO2 emissions}\label{how-gcam-represents-non-co2-emissions}}

A non-CO2
emission
in GCAM
exists
only where
the model
attaches
an
emissions
``tag'' to
a
technology
(audit
§A1). The
output
database
stores
just the
emitted
mass per
year. The
emission
factor is
not
stored; we
compute it
as mass ÷
activity.

The mass
gets there
in two
phases
(audit
§A2):

\begin{itemize}
\tightlist
\item
  \textbf{Base
  years
  (up to
  2021).}
  An
  inventory
  total is
  supplied,
  and GCAM
  divides
  it by
  its own
  activity
  to get
  the
  factor.
\item
  \textbf{Future
  years.}
  That
  factor
  is
  carried
  forward,
  either
  frozen
  or
  reduced
  along a
  pre-set
  curve,
  and
  multiplied
  by the
  scenario's
  activity.
\end{itemize}

\hypertarget{the-investigation-step-by-step}{%
\subsection*{3. The investigation, step by step}\addcontentsline{toc}{subsection}{3. The investigation, step by step}\label{the-investigation-step-by-step}}

Each step
is laid
out as:
the
question
we asked,
what we
did, what
we found,
and a
plain-English
takeaway.

\hypertarget{step-1.-compare-gcam-kaist-with-capss-by-source-category}{%
\subsubsection*{Step 1. Compare GCAM-KAIST with CAPSS by source category}\addcontentsline{toc}{subsubsection}{Step 1. Compare GCAM-KAIST with CAPSS by source category}\label{step-1.-compare-gcam-kaist-with-capss-by-source-category}}

\textbf{Question.}
How close
are
GCAM-KAIST's
base-year
emissions
to CAPSS?

\textbf{What
we did.}

\begin{itemize}
\tightlist
\item
  Assigned
  each of
  127
  GCAM-KAIST
  emitting
  sectors
  to one
  of
  CAPSS's
  13
  first-level
  source
  categories
  (crosswalk:
  \texttt{src/\allowbreak{}gcam\_\allowbreak{}emissions/\allowbreak{}reference/\allowbreak{}gcam\_\allowbreak{}kaist\_\allowbreak{}capss\_\allowbreak{}category\_\allowbreak{}crosswalk.\allowbreak{}csv}).
\item
  Applied
  two
  finer
  overrides:
  residue-burning
  species
  →
  Biomass
  burning,
  and the
  \texttt{solvents}
  subsector
  →
  Solvent
  use.
\item
  Excluded
  international
  aviation
  and
  shipping
  as
  bunker
  fuel.
\item
  Compared
  the
  totals
  with
  CAPSS's
  published
  2015 and
  2021
  category
  totals
  for the
  six
  pollutants
  both
  report
  (SOx,
  NOx,
  VOCs,
  NH3, CO,
  BC).
  GCAM
  reports
  no PM2.5
  or PM10.
\end{itemize}

\textbf{What
we found.}
Agreement
is very
uneven
(Figure
1).

\begin{itemize}
\tightlist
\item
  Some
  cells
  are
  nearly
  identical:
  agricultural
  NH3,
  road-transport
  NOx in
  2021,
  industrial-process
  NH3,
  waste
  SOx.
\item
  Others
  are far
  off:

  \begin{itemize}
  \tightlist
  \item
    \textbf{BC}
    is
    9--65×
    higher
    than
    CAPSS
    in the
    stationary-combustion
    categories.
  \item
    \textbf{Biomass-burning
    CO and
    VOCs}
    are
    only
    6--11\%
    of
    CAPSS.
  \item
    \textbf{Energy
    transport
    and
    storage
    VOCs}
    (fuel
    storage
    and
    distribution)
    are
    only
    12--18\%
    of
    CAPSS.
  \end{itemize}
\end{itemize}

\figslot{fig01_capss_vs_gcam_2021.png}

\emph{Figure
1. CAPSS
vs
GCAM-KAIST
by CAPSS
category,
2021. The
2015
version is
\texttt{figures/\allowbreak{}fig01\_\allowbreak{}capss\_\allowbreak{}vs\_\allowbreak{}gcam\_\allowbreak{}2015.\allowbreak{}png}.}

\begin{quote}
\plainmark{}
\textbf{Takeaway.}
Some
categories
match
almost
perfectly
and others
are off by
an order
of
magnitude.
Before
reading
the good
matches as
``GCAM is
accurate'',
we needed
to know
where
GCAM's
numbers
come from.
\end{quote}

\hypertarget{step-2.-understand-how-gcam-stores-and-reports-non-co2}{%
\subsubsection*{Step 2. Understand how GCAM stores and reports non-CO2}\addcontentsline{toc}{subsubsection}{Step 2. Understand how GCAM stores and reports non-CO2}\label{step-2.-understand-how-gcam-stores-and-reports-non-co2}}

\textbf{Question.}
Is there a
GCAM
output
we're
missing,
such as a
query for
cement or
refinery
emissions?

\textbf{What
we did.}
Read the
output
database
structure
and GCAM's
standard
reporting
queries
(audit
§B2, §B3).

\textbf{What
we found.}

\begin{itemize}
\tightlist
\item
  GCAM has
  no
  sector-specific
  non-CO2
  queries.
  The
  standard
  queries
  return
  every
  tag that
  exists.
\item
  A sector
  reports
  non-CO2
  emissions
  if, and
  only if,
  its
  technologies
  carry
  tags.
\end{itemize}

\begin{quote}
\plainmark{}
\textbf{Takeaway.}
If
GCAM-KAIST
reports
nothing
for an
activity,
it isn't a
missing
query. The
model
never
attached
emissions
to that
activity.
\end{quote}

\hypertarget{step-3.-check-which-activities-carry-emissions-tags}{%
\subsubsection*{Step 3. Check which activities carry emissions tags}\addcontentsline{toc}{subsubsection}{Step 3. Check which activities carry emissions tags}\label{step-3.-check-which-activities-carry-emissions-tags}}

\textbf{Question.}
Does every
emitting
activity
in
GCAM-KAIST
have tags?

\textbf{What
we did.}
Listed
every
technology
with
activity
in either
scenario
up to
2050, and
checked
which
carry tags
(audit
§B4).

\textbf{What
we found.}

\begin{itemize}
\tightlist
\item
  All
  operating
  fossil
  power
  plants
  are
  tagged.
\item
  Several
  large
  activities
  are not:

  \begin{itemize}
  \tightlist
  \item
    the
    cement
    \emph{process}
    (only
    kiln
    fuel
    is
    tagged);
  \item
    oil
    refining
    (4.69
    EJ of
    output
    in
    2021);
  \item
    chemical
    feedstocks
    and
    industrial
    cogeneration.
  \end{itemize}
\item
  Several
  technologies
  that
  grow in
  the
  net-zero
  scenario
  are also
  untagged:
  hybrid
  vehicles,
  hydrogen-based
  steel,
  and
  steel
  and
  cement
  with
  carbon
  capture.
\item
  All of
  these
  gaps
  also
  exist in
  stock
  GCAM v9
  (audit
  §B4).
\end{itemize}

\begin{quote}
\plainmark{}
\textbf{Takeaway.}
Some real
sources of
air
pollution
are simply
absent
from GCAM,
and the
net-zero
scenario
moves
activity
toward
technologies
that the
model
treats as
emission-free.
\end{quote}

\hypertarget{step-4.-trace-where-the-base-year-numbers-come-from}{%
\subsubsection*{Step 4. Trace where the base-year numbers come from}\addcontentsline{toc}{subsubsection}{Step 4. Trace where the base-year numbers come from}\label{step-4.-trace-where-the-base-year-numbers-come-from}}

\textbf{Question.}
Why do
some CAPSS
categories
match so
closely?

\textbf{What
we did.}

\begin{itemize}
\tightlist
\item
  Compared
  GCAM-KAIST's
  base-year
  tags
  with
  stock
  GCAM
  v9's
  input
  files
  (audit
  §B5.2).
\item
  Traced
  stock
  GCAM's
  calibration
  source
  through
  CEDS's
  documentation
  (audit
  §B5.3).
\item
  Measured
  how
  closely
  each
  category
  ×
  pollutant
  × year
  cell
  agrees
  with
  CAPSS
  (audit
  §B5.4).
\end{itemize}

\textbf{What
we found.}

\begin{itemize}
\tightlist
\item
  GCAM-KAIST's
  base-year
  emissions
  are
  stock
  GCAM
  v9's:
  2,244 of
  2,251
  values
  match
  within
  0.1\%.
\item
  Stock GCAM calibrates to CEDS, and CEDS's own code rescales its Korean
  estimates to CAPSS, sector group by sector group (audit §A3 item 2 walks
  through the steps; audit §B5.3 has the code evidence).
\item
  For 2000--2018, CEDS then rescales Korea a second time, to the EDGAR-HTAP
  inventory. 2021 is outside that window, which makes it the cleaner year for
  comparing with CAPSS.
\item
  \textbf{11
  of the
  102
  comparable
  cells
  agree
  with
  CAPSS to
  within
  0.12\%},
  several
  to four
  or five
  significant
  figures.
  For
  example,
  agricultural
  NH3 is
  200,384.0
  t vs
  CAPSS
  200,384
  t in
  2021
  (Figure
  2).
\end{itemize}

\figslot{fig02_capss_agreement.png}

\emph{Figure 2. Relative difference from CAPSS for every comparable cell
(log scale). Each dot is one pair of bars from Figure 1 (2021) or its 2015
version: the blue CAPSS bar and the orange GCAM-KAIST bar for the same
category and pollutant. Its position is the percent difference between the
two, $|$GCAM-KAIST / CAPSS $-$ 1$|$. A cell is comparable when both bars are
non-zero, which gives 102 cells across the two years. Blue has a different
meaning here than in Figure 1: the 11 blue dots agree within 0.12\%, and the
rest differ by 1\% to over 1,000\%.}

\textbf{Why do only 11 cells match?} Among the other 91 cells, 13 are within
5\% of CAPSS and 71 are more than 10\% off. The calibration chain (CAPSS →
CEDS → stock GCAM → GCAM-KAIST) explains both patterns. It forces agreement
with CAPSS only at the level where CEDS rescales to it, and each step after
that can move mass across CAPSS's category boundaries. Our proposed
explanations are below. The first four rest on evidence in this document or
the audit; the fifth is untested.

\begin{enumerate}
\def\labelenumi{\arabic{enumi}.}
\item
  \textbf{A CAPSS total survives only where the categories line up at every
  step.}

  \begin{itemize}
  \tightlist
  \item
    CEDS rescales its Korean estimates to CAPSS one sector group at a time.
  \item
    Inside each group, CEDS's own sector shares divide the total.
  \item
    GCAM then maps CEDS sectors onto its own sectors, and we map those onto
    CAPSS categories.
  \item
    A CAPSS category total comes back out intact only if all of these
    boundaries cover the same sources.
  \item
    All 11 matching cells fall in four places: agricultural NH3, road
    transport, industrial-process NH3, and waste SOx. These appear to be the
    places where the boundaries coincide (inferred).
  \item
    Elsewhere, mass that CAPSS counts in one category lands in another.
    Under our crosswalk, 44.4\% of CAPSS's emitted mass ends up in the wrong
    category (Step 5). Step 5 also shows that no reassignment of GCAM sectors
    repairs this, so the boundaries are lost before our crosswalk is applied.
  \end{itemize}
\item
  \textbf{Some pollutants were never rescaled to CAPSS.} For Korean BC, CEDS
  rescales only road transport (audit §B5.3). Everywhere else, BC keeps
  CEDS's own estimate. That fits the pattern: road-transport BC is within 5\%
  of CAPSS, while stationary-combustion BC is 9--65$\times$ CAPSS (audit
  §B5.4).
\item
  \textbf{A second rescaling can pull 2015 away from CAPSS.} For 2000--2018,
  CEDS also rescales Korea to EDGAR-HTAP (audit §B5.3).

  \begin{itemize}
  \tightlist
  \item
    Road NOx agrees with CAPSS within 0.03\% in 2021, which is outside that
    window, but is 25\% high in 2015.
  \item
    The effect is uneven: five of the 11 matching cells are from 2015.
  \item
    We haven't checked which Korean sectors this step actually changes.
  \end{itemize}
\item
  \textbf{Some sources sit in a different place in GCAM, or nowhere.}

  \begin{itemize}
  \tightlist
  \item
    Refinery fuel combustion appears to sit in GCAM's industrial energy use
    rather than with energy production (audit §B5.5, inferred).
  \item
    Several activities carry no tag at all (Step 3).
  \item
    Two categories are mostly missing:

    \begin{itemize}
    \tightlist
    \item
      \textbf{Energy transport and storage VOCs} are only 12--18\% of CAPSS.
      CAPSS's category covers fuel storage and distribution losses, but the
      GCAM rows we map there are resource-extraction fugitives and hydrogen
      delivery. We found no GCAM tag that covers fuel distribution.
    \item
      \textbf{Biomass-burning CO and VOCs} are only 6--11\% of CAPSS. In
      GCAM, this category holds only crop-residue burning and forest,
      grassland and deforestation fires; residential wood burning is mapped
      to Non-industry. The remainder of CAPSS's category is likely missing or
      sitting in another category.
    \end{itemize}
  \end{itemize}
\item
  \textbf{Different CAPSS editions (untested).} CAPSS revises past years. CEDS
  keeps its own copy of the inventory
  (\texttt{Korea\_CAPSS\_Emissions.xlsx}), which may be an older edition than
  the totals we compared against. That would produce misses even where the
  categories line up. The 11 exact matches show that, at least for those
  cells, the editions agree.
\end{enumerate}

\begin{quote}
\plainmark{}
\textbf{Takeaway.}
Agreement
to five
significant
figures
doesn't
happen by
independent
estimation.
These
base-year
numbers
were built
from the
inventory
itself, so
their
agreement
with CAPSS
is
circular:
it shows
the
calibration
worked,
not that
GCAM
represents
emissions
correctly. The 91
misses aren't independent evidence either. They mark the places where the
chain blurred CAPSS's categories, never rescaled a pollutant to CAPSS, or
dropped a source.
\end{quote}

\hypertarget{step-5.-stress-test-the-category-mapping}{%
\subsubsection*{Step 5. Stress-test the category mapping}\addcontentsline{toc}{subsubsection}{Step 5. Stress-test the category mapping}\label{step-5.-stress-test-the-category-mapping}}

\textbf{Question.}
Is the
poor
agreement
elsewhere
just a
mapping
problem?
Would a
better
assignment
of GCAM
sectors to
CAPSS
categories
make
everything
line up?

\textbf{What
we did.}

\begin{itemize}
\tightlist
\item
  Identified
  15
  groups
  of
  GCAM-KAIST
  rows
  whose
  CAPSS
  category
  is
  ambiguous.
  Examples:
  stationary
  construction
  machinery,
  residential
  wood
  burning,
  iron and
  steel,
  international
  shipping
  and
  aviation,
  and
  industrial
  oil use.
\item
  Gave
  each
  group
  its 2--3
  plausible
  categories,
  and
  evaluated
  all
  49,152
  combinations
  against
  CAPSS
  for 2015
  and
  2021.
\item
  Scored
  each
  mapping
  by the
  share of
  CAPSS's
  emitted
  mass
  that
  ends up
  in the
  wrong
  category,
  averaged
  over the
  12 year
  ×
  pollutant
  combinations.
\end{itemize}

Script:
\texttt{results/\allowbreak{}diagnostics/\allowbreak{}capss\_\allowbreak{}calibration\_\allowbreak{}test/\allowbreak{}mapping\_\allowbreak{}search.\allowbreak{}py}.

\textbf{What
we found.}

\begin{longtable}[]{@{}
  >{\raggedright\arraybackslash}p{(\columnwidth - 8\tabcolsep) * \real{0.1855}}
  >{\raggedright\arraybackslash}p{(\columnwidth - 8\tabcolsep) * \real{0.3710}}
  >{\raggedright\arraybackslash}p{(\columnwidth - 8\tabcolsep) * \real{0.1935}}
  >{\raggedright\arraybackslash}p{(\columnwidth - 8\tabcolsep) * \real{0.1210}}
  >{\raggedright\arraybackslash}p{(\columnwidth - 8\tabcolsep) * \real{0.1290}}@{}}
\toprule\noalign{}
\begin{minipage}[b]{\linewidth}\raggedright
Mapping
\end{minipage}
&
\begin{minipage}[b]{\linewidth}\raggedright
Share of
CAPSS mass
in the
wrong
category
\end{minipage}
&
\begin{minipage}[b]{\linewidth}\raggedright
Cells
within
0.12\%
\end{minipage}
&
\begin{minipage}[b]{\linewidth}\raggedright
Within 5\%
\end{minipage}
&
\begin{minipage}[b]{\linewidth}\raggedright
Within
10\%
\end{minipage} \\
\midrule\noalign{}
\endhead
\bottomrule\noalign{}
\endlastfoot
Current
crosswalk
& 44.4\% &
11 & 24 &
31 \\
Best of
49,152 &
40.9\% &
12 & 29 &
41 \\
\end{longtable}

\emph{(out
of 102
comparable
cells)}

\begin{itemize}
\tightlist
\item
  No
  mapping
  comes
  close to
  reproducing
  CAPSS;
  the best
  removes
  only 3.5
  percentage
  points
  of
  misallocation.
  The
  near-exact
  cells in
  Step 4
  appear
  under
  the
  current
  crosswalk,
  with no
  search
  needed.
\item
  The best
  mapping
  makes
  four
  changes:

  \begin{itemize}
  \tightlist
  \item
    residential
    traditional
    biomass
    →
    Biomass
    burning;
  \item
    biodiesel
    refining
    →
    Manufacturing
    industry;
  \item
    international
    aviation
    →
    Non-road
    transport;
  \item
    \texttt{other\ industrial\ energy\ use\ /\allowbreak{}\ refined\ liquids}
    →
    Energy
    production.
  \end{itemize}
\item
  The last
  one is
  informative.
  It
  brings
  Energy
  production
  CO and
  NH3
  within
  2\% of
  CAPSS.
  Together
  with
  GCAM's
  own CEDS
  sector
  map, it
  indicates
  that
  refinery
  fuel-combustion
  emissions
  sit
  inside
  GCAM's
  industrial
  energy
  use
  (audit
  §B5.5,
  an
  inference).
\end{itemize}

\begin{quote}
\plainmark{}
\textbf{Takeaway.}
The
mismatches
are not a
bookkeeping
problem we
can fix by
relabelling.
Agreement
is
excellent
where a
GCAM
sector
lines up
one-to-one
with a
CAPSS
category,
and poor
where it
doesn't.
That
pattern is
what
calibration
produces,
not what
physics
produces.
\end{quote}

\emph{Caveat.}
With 15
binary or
ternary
choices,
some
improvement
is
expected
by chance.
The search
is used
only to
show that
remapping
can't
\emph{create}
agreement;
the
near-exact
cells need
no search.

\hypertarget{step-6.-back-calculate-the-emission-factors}{%
\subsubsection*{Step 6. Back-calculate the emission factors}\addcontentsline{toc}{subsubsection}{Step 6. Back-calculate the emission factors}\label{step-6.-back-calculate-the-emission-factors}}

\textbf{Question.}
What
emission
factor
does each
GCAM-KAIST
technology
implicitly
use, and
does it
reflect
the
technology?

\textbf{What
we did.}
Divided
each tag's
emissions
by the
activity
GCAM ties
it to
(fuel
burned, or
output)
for every
base year
(audit
§B6.1).

\textbf{What
we found.}

\textbf{(a)
Different
technologies
share one
factor.}

\begin{itemize}
\tightlist
\item
  Combined-cycle
  and
  steam/turbine
  gas
  plants
  have
  identical
  NOx and
  CO
  factors
  per unit
  of fuel
  in every
  base
  year
  from
  2005 to
  2021
  (Figure
  3).
\item
  Mobile
  and
  stationary
  machinery
  carry
  exactly
  equal
  emission
  masses
  in every
  base
  year
  (Figure
  4).
\item
  Commercial
  heating,
  cooling
  and
  ``other''
  uses of
  gas
  share
  one
  factor,
  and so
  do all
  ten
  household
  income
  groups
  (audit
  §B6.2).
\end{itemize}

\figslot{fig03_gas_cc_vs_steamct.png}

\emph{Figure
3. Implied
NOx and CO
factors
per TJ of
gas for
the two
gas-power
technologies.
They
coincide
in every
base year
and
diverge
only after
2021.}

\figslot{fig04_mobile_vs_stationary.png}

\emph{Figure
4. NOx
from
refined-liquids
machinery:
mobile and
stationary
carry
identical
masses in
every base
year, in
both
agriculture
and
construction.
(Mining,
not shown,
behaves
the
same.)}

\textbf{(b)
Factors
swing
between
base
years.}
The
typical
technology-pollutant
factor
varies
about 2×
across
2005--2021,
and one in
ten varies
more than
11×
(Figure
5). A real
emission
rate
doesn't
move like
that.

\figslot{fig05_base_year_swing.png}

\emph{Figure
5. Ratio
of the
largest to
the
smallest
implied
factor
across
2005,
2010, 2015
and 2021,
for 436
technology-pollutant
series.}

\textbf{(c)
Large
shares of
emissions
are tied
to no
activity
at all.}
In 2021,
81\% of
GCAM-KAIST's
Korean
VOCs and
55\% of
its SOx
sit on
tags
driven by
a fixed
placeholder,
not by any
production
quantity
(Figure 6;
audit
§B6.4).
This
excludes
international
shipping
and
aviation.

\figslot{fig06_placeholder_shares.png}

\emph{Figure
6. Share
of 2021
emissions
on
placeholder-driven
tags
(solvents,
other
industrial
processes,
waste).}

\begin{quote}
\plainmark{}
\textbf{Takeaway.}
GCAM's
factors
are what's
left after
dividing
an
inventory
total
among
technologies.
They can't
tell a
combined-cycle
gas plant
from an
older one,
or a
tractor
from a
farm
boiler.
They jump
around
between
years, and
for most
VOCs and
SOx they
aren't
attached
to
anything
that
changes
with
production.
\end{quote}

\hypertarget{step-7.-follow-the-factors-into-the-future}{%
\subsubsection*{Step 7. Follow the factors into the future}\addcontentsline{toc}{subsubsection}{Step 7. Follow the factors into the future}\label{step-7.-follow-the-factors-into-the-future}}

\textbf{Question.}
How do the
factors
behave in
the
scenario
years we
actually
need?

\textbf{What
we did.}
Compared
each
factor in
2050 with
its 2021
value, and
the
\texttt{nz}
value with
the
\texttt{ref}
value
(audit
§B6.5).

\textbf{What
we found.}

\begin{itemize}
\tightlist
\item
  \textbf{After
  2021,
  46\% of
  factors
  stay
  frozen
  and 52\%
  decline
  along
  the
  pre-set
  curve.}
\item
  \textbf{The pre-set curve is GCAM's ``GDP control''.} It assumes pollution
  controls tighten as a country gets richer.

  \begin{itemize}
  \tightlist
  \item
    After the last base year (2021), GCAM multiplies each controlled factor
    by 1 / (1 + (GDP per capita in that year $-$ GDP per capita in 2021) /
    steepness).
  \item
    Each factor's \texttt{max-reduction} sets the floor. A factor with
    \texttt{max-reduction} 60, for example, never drops below 40\% of its 2021
    value.
  \item
    Source: \texttt{GDPControl::calcEmissionsReduction} in the GCAM v9.0
    source code
    (\texttt{cvs/\allowbreak{}objects/\allowbreak{}emissions/\allowbreak{}source/\allowbreak{}gdp\_\allowbreak{}control.\allowbreak{}cpp})
    and GCAM's emissions documentation.
  \item
    In stock GCAM v9, South Korea has 469 of these controls. All of them use
    steepness 3.5, and their \texttt{max-reduction} values range from 0 to
    100. A factor stays frozen when it has no GDP control, or when its
    \texttt{max-reduction} is 0.
  \item
    GCAM-KAIST's own control settings can't be inspected, because its input
    files aren't available. Its outputs, however, show that one curve does
    almost all the work. Of the 243 declining factors, 195 fall to exactly
    0.515× their 2021 value by 2030, and 161 fall to 0.289× by 2050. The
    decline is the same for coal power, cars, chemical-industry boilers and
    commercial buildings, whatever the pollutant. Of the other 82, 68 look
    like the same curve stopped at a \texttt{max-reduction} floor: they either
    stay flat from 2030 to 2050, or follow 0.515× in 2030 and then level off
    above 0.289× (for example, 0.4× in 2050 matches a \texttt{max-reduction}
    of 60). The remaining 14 include international shipping SO2, which
    \texttt{emission\_factor\_controls.xml} sets directly (audit §B5.1), and
    several gas power and natural-gas car factors.
  \item
    The curve depends only on GDP per capita. The net-zero scenario evidently
    assumes the same GDP path, because the shared multipliers are identical
    in \texttt{ref} and \texttt{nz}.
  \end{itemize}
\item
  \textbf{In 2050, 93\% are
  identical
  in the
  reference
  and
  net-zero
  scenarios}
  (Figure
  7).
\item
  The
  net-zero
  scenario
  therefore
  changes
  emissions
  almost
  only
  through
  activity:
  how much
  of each
  technology
  runs. It
  does not
  change
  how
  clean
  each
  technology
  is.
\end{itemize}

\figslot{fig07_coefficient_trajectories.png}

\emph{Figure
7. Four
examples.
The
\texttt{nz}
markers
sit
exactly on
the
\texttt{ref}
line. Coal
power and
car
factors
decline
along the
pre-set
curve;
cement-kiln
coal and
commercial
gas
factors
stay
frozen at
their 2021
values.}

\begin{quote}
\plainmark{}
\textbf{Takeaway.}
The model
has no
lever for
pollution-control
policy:
scrubbers,
emission
standards,
cleaner
boilers.
Any such
improvement
is assumed
identically
in both
scenarios,
so it can
never show
up as a
benefit of
the
policy.
\end{quote}

\hypertarget{does-it-change-the-policy-answer-the-policy-experiment}{%
\subsection*{4. Does it change the policy answer? The policy experiment}\addcontentsline{toc}{subsection}{4. Does it change the policy answer? The policy experiment}\label{does-it-change-the-policy-answer-the-policy-experiment}}

\hypertarget{what-the-experiment-does}{%
\subsubsection*{4.1 What the experiment does \plainmark{}}\addcontentsline{toc}{subsubsection}{4.1 What the experiment does}\label{what-the-experiment-does}}

Emissions
=
\textbf{activity
× emission
factor},
technology
by
technology.
GCAM-KAIST
gives us
both.

The
experiment
keeps
\textbf{GCAM-KAIST's
activity
exactly as
it is} in
both
scenarios
(how much
coal power
runs, how
much steel
each route
makes, how
far hybrid
buses
travel)
and
\textbf{swaps
only the
emission
factors}.
Then it
asks the
policy
question:

\begin{quote}
\emph{How
much lower
are 2050
emissions
in the
net-zero
scenario
than in
the
reference
scenario?}
\end{quote}

If GCAM's
own
factors
and more
realistic
factors
give about
the same
answer,
GCAM's
native
emissions
would be
good
enough for
comparing
policies,
even if
their
levels
were off.
If the
answers
differ,
the native
emissions
can't be
used to
judge the
policy.

We ran it
for three
sectors,
each
testing a
different
weakness
from
Section 3:

\begin{itemize}
\tightlist
\item
  \textbf{Power}
  tests
  ``one
  factor
  per
  fuel, on
  a curve
  shared
  by both
  scenarios''
  (Steps
  6--7).
\item
  \textbf{Steel}
  tests
  ``untagged
  technologies
  count as
  zero''
  (Step
  3).
\item
  \textbf{Road
  transport}
  tests
  the same
  thing
  for
  hybrid
  vehicles
  (Step
  3).
\end{itemize}

\hypertarget{design}{%
\subsubsection*{4.2 Design}\addcontentsline{toc}{subsubsection}{4.2 Design}\label{design}}

\begin{longtable}[]{@{}
  >{\raggedright\arraybackslash}p{(\columnwidth - 4\tabcolsep) * \real{0.1500}}
  >{\raggedright\arraybackslash}p{(\columnwidth - 4\tabcolsep) * \real{0.2750}}
  >{\raggedright\arraybackslash}p{(\columnwidth - 4\tabcolsep) * \real{0.5750}}@{}}
\toprule\noalign{}
\begin{minipage}[b]{\linewidth}\raggedright
Sector
\end{minipage}
&
\begin{minipage}[b]{\linewidth}\raggedright
Variant
\end{minipage}
&
\begin{minipage}[b]{\linewidth}\raggedright
Emission
factors
used
\end{minipage} \\
\midrule\noalign{}
\endhead
\bottomrule\noalign{}
\endlastfoot
Power &
native &
GCAM-KAIST's
own
emissions \\
& GCAM
2021,
fixed &
GCAM's own
2021
factor per
MWh for
each
technology,
held
constant
to 2050 \\
& KEPCO 2021 & Our hand-calculated KEPCO 2021 factors per MWh × GCAM
generation, by technology (kg/MWh): coal (NOx 0.126, SOx 0.139); gas (NOx
0.164, SOx 0); oil (NOx 0.301, SOx 0.011); biomass (NOx 0.180, SOx 0.003) \\
Steel &
native &
GCAM-KAIST's
own
emissions;
hydrogen-based
DRI and
blast
furnace
with CCS
count as
zero \\
& filled &
Hydrogen-based
DRI at the
CAPSS
electric-arc-furnace
process
factor
(NOx 0.2,
SOx 0.35,
VOCs 0.09
kg/t);
blast
furnace
with CCS
at GCAM's
own
blast-furnace
factor \\
Road &
native &
GCAM-KAIST's
own
emissions;
hybrids
count as
zero \\
& filled &
Hybrids at
the same
vehicle
type's
conventional-fuel
factor per
unit of
fuel \\
\end{longtable}

\textbf{What the power comparison covers.}

\begin{itemize}
\tightlist
\item
  \textbf{Sector.} GCAM-KAIST's \texttt{electricity} sector: the
  technologies in its coal, gas, refined liquids and biomass subsectors.
  Nuclear and renewables have no tags in any variant, and the
  \texttt{elec\_*} cooling-system pass-through sectors carry no emissions.
  Industrial cogeneration is not included.
\item
  \textbf{Quantity.} NOx and SOx in 2050, reference minus net-zero, in kt.
  Each variant applies the same factors to both scenarios.
\item
  \textbf{Activity.} GCAM-KAIST's electricity output by technology (EJ,
  converted to MWh). In 2050 the reference scenario generates 101.0 TWh from
  coal, 99.6 TWh from gas combined cycle, 3.9 TWh from gas steam/turbine,
  5.1 TWh from biomass and 0.25 TWh from oil. The net-zero scenario
  generates 2.9 TWh from gas combined cycle with CCS and almost nothing else
  from these fuels.
\end{itemize}

\textbf{Where the KEPCO factors come from.} We calculated them by hand in
NZK-APHIAM from the monthly generation and stack-emission records of
KEPCO's five thermal power subsidiaries (East-West, Western, Southern,
South-East and Midland).

\begin{itemize}
\tightlist
\item
  Each factor is the generation-weighted average of plant-level factors (kg
  emitted per MWh generated) for one fuel × technology cohort in 2021, under
  NZK-APHIAM's \texttt{operational\_primary} screening rule.
\item
  Table:
  \texttt{NZK-APHIAM/\allowbreak{}data/\allowbreak{}processed/\allowbreak{}kepco/\allowbreak{}emission\_\allowbreak{}factors/\allowbreak{}kepco\_\allowbreak{}annual\_\allowbreak{}ef\_\allowbreak{}distribution\_\allowbreak{}long\_\allowbreak{}by\_\allowbreak{}fuel\_\allowbreak{}technology.\allowbreak{}csv}.
\item
  Plausibility screen:
  \texttt{NZK-APHIAM/\allowbreak{}docs/\allowbreak{}references/\allowbreak{}thermal/\allowbreak{}kepco\_\allowbreak{}ef\_\allowbreak{}external\_\allowbreak{}validation.\allowbreak{}md}.
\end{itemize}

\begin{longtable}[]{@{}p{0.22\linewidth}p{0.27\linewidth}p{0.12\linewidth}p{0.10\linewidth}p{0.10\linewidth}@{}}
\toprule\noalign{}
GCAM technology & KEPCO 2021 cohort & Plants & NOx (kg/\allowbreak{}MWh) & SOx (kg/\allowbreak{}MWh) \\
\midrule\noalign{}
\endhead
\bottomrule\noalign{}
\endlastfoot
\texttt{coal\ (conv\ pul)} & coal, conventional steam turbine & 9 & 0.126 & 0.139 \\
\texttt{gas\ (CC)}, \texttt{gas\ (CC\ CCS)} & natural gas, combined cycle & 8 (NOx); 6 (SOx) & 0.164 & 0 (reported) \\
\texttt{gas\ (steam/CT)} & same combined-cycle cohort; KEPCO has no 2021 gas steam/turbine cohort & & 0.164 & 0 \\
\texttt{refined\ liquids\ (steam/CT)} & oil, conventional steam turbine & 1 & 0.301 & 0.011 \\
\texttt{biomass\ (conv)}, \texttt{biomass\ (IGCC)} & biomass, conventional steam turbine & 1 & 0.180 & 0.003 \\
\end{longtable}

Notes on the design:

\begin{itemize}
\tightlist
\item
  The
  steel
  fill is
  a
  \textbf{lower
  bound}:
  hydrogen-based
  DRI also
  burns
  gas and
  oil
  (audit
  §B4),
  and that
  combustion
  isn't
  counted.
\end{itemize}

Script:
\texttt{results/\allowbreak{}diagnostics/\allowbreak{}policy\_\allowbreak{}experiment/\allowbreak{}policy\_\allowbreak{}experiment.\allowbreak{}py}.

\hypertarget{results}{%
\subsubsection*{4.3 Results}\addcontentsline{toc}{subsubsection}{4.3 Results}\label{results}}

\textbf{Power:
same
direction,
very
different
size}
(Figure
8). The
net-zero
scenario
has almost
no fossil
power left
in 2050,
so every
variant
shows a
cut of
about
99\%. What
differs is
how much
is
\emph{avoided},
because
that
depends on
how dirty
the
reference
scenario's
power is:

\begin{longtable}[]{@{}
  >{\raggedright\arraybackslash}p{(\columnwidth - 4\tabcolsep) * \real{0.4359}}
  >{\raggedright\arraybackslash}p{(\columnwidth - 4\tabcolsep) * \real{0.2821}}
  >{\raggedright\arraybackslash}p{(\columnwidth - 4\tabcolsep) * \real{0.2821}}@{}}
\toprule\noalign{}
\begin{minipage}[b]{\linewidth}\raggedright
2050,
power
sector
\end{minipage}
&
\begin{minipage}[b]{\linewidth}\raggedright
NOx
avoided
(kt)
\end{minipage}
&
\begin{minipage}[b]{\linewidth}\raggedright
SOx
avoided
(kt)
\end{minipage} \\
\midrule\noalign{}
\endhead
\bottomrule\noalign{}
\endlastfoot
GCAM-KAIST
native &
7.5 &
2.7 \\
GCAM's
2021
factors,
fixed &
19.7
(2.6×) &
7.4
(2.7×) \\
KEPCO 2021 factors & 30.3 (4.0×) & 14.1 (5.2×) \\
\end{longtable}

Holding
GCAM's
\emph{own}
factors
fixed
already
gives
2.6×. Most of the gap between native and KEPCO therefore comes from the pre-set curve that
shrinks
the
reference
scenario's
factors.
The curve
is an
assumption
shared by
both
scenarios,
not a
policy
outcome.

For context, in 2021 KEPCO's 2021 factors × GCAM generation give 1.5× GCAM's native power NOx and 1.8× its SOx.

\figslot{fig08_policy_power.png}

\emph{Figure
8. NOx and
SOx
avoided by
the
net-zero
scenario
in 2050,
power
sector,
under each
set of
factors.}

\textbf{Steel:
the
direction
flips}
(Figure
9).

\begin{longtable}[]{@{}
  >{\raggedright\arraybackslash}p{(\columnwidth - 6\tabcolsep) * \real{0.5833}}
  >{\raggedright\arraybackslash}p{(\columnwidth - 6\tabcolsep) * \real{0.1389}}
  >{\raggedright\arraybackslash}p{(\columnwidth - 6\tabcolsep) * \real{0.1389}}
  >{\raggedright\arraybackslash}p{(\columnwidth - 6\tabcolsep) * \real{0.1389}}@{}}
\toprule\noalign{}
\begin{minipage}[b]{\linewidth}\raggedright
2050,
steel
sector, nz
vs ref
\end{minipage}
&
\begin{minipage}[b]{\linewidth}\raggedright
NOx
\end{minipage}
&
\begin{minipage}[b]{\linewidth}\raggedright
SOx
\end{minipage}
&
\begin{minipage}[b]{\linewidth}\raggedright
VOCs
\end{minipage} \\
\midrule\noalign{}
\endhead
\bottomrule\noalign{}
\endlastfoot
GCAM-KAIST
native &
−35\% &
\textbf{−88\%}
& −53\% \\
Untagged
routes
filled
(lower
bound) &
−27\% &
\textbf{+19\%}
&
\textbf{+26\%} \\
\end{longtable}

The 88\%
SOx cut
comes from
the
net-zero
scenario's
steel
moving
onto
routes
with no
tags (16.4
Mt
hydrogen-based
DRI, 4.2
Mt blast
furnace
with CCS).
Once their
electric-arc-furnace
process
emissions
are
counted,
steel SOx
and VOCs
\emph{rise}
under net
zero.

\figslot{fig09_policy_steel.png}

\emph{Figure
9. Change
in 2050
steel-sector
emissions,
net zero
vs
reference.}

\textbf{Road:
modest}
(Figure
10).
Counting
hybrids
raises
2050 road
emission
levels in
both
scenarios:

\begin{itemize}
\tightlist
\item
  NOx:
  22--26\%
  higher;
\item
  CO:
  40--47\%
  higher;
\item
  VOCs:
  41--48\%
  higher.
\end{itemize}

The
relative
cut barely
changes
(NOx −48\%
native vs
−46\%
filled).

\figslot{fig10_policy_road.png}

\emph{Figure
10. 2050
road
emissions
with and
without
hybrids
counted.}

\begin{quote}
\plainmark{}
\textbf{Takeaway.}
Using
GCAM-KAIST's
native
emissions
would get
the
\emph{direction}
wrong for
steel, and
the
\emph{size}
of the
power-sector
benefit too small by 2.6--5.2×.
Road
emission
levels
would be
too low by
a fifth to
a third or
more.
Because
the health
benefit of
a policy
is
computed
from the
emissions
it avoids,
these
errors
carry
straight
through to
the health
results.
\end{quote}

\hypertarget{limits-of-the-experiment}{%
\subsubsection*{4.4 Limits of the experiment}\addcontentsline{toc}{subsubsection}{4.4 Limits of the experiment}\label{limits-of-the-experiment}}

\begin{itemize}
\tightlist
\item
  \textbf{It is illustrative.} None of the KEPCO or CAPSS factors used are
  production-ready. The steel factor is from CAPSS VI and is flagged as
  superseded pending the CAPSS VII update.
\item
  \textbf{The KEPCO factors cover KEPCO's plants, not every Korean plant.}
  They come from KEPCO's five thermal subsidiaries, which excludes independent
  power producers. The oil and biomass cohorts are one plant each. Gas SOx is
  reported as zero for a cohort covering 57\% of its generation.
\item
  \textbf{The KEPCO factors are held at 2021 values.} NZK-APHIAM also has
  later years (2025 coal NOx is 0.087 kg/MWh, for example). Projecting those
  factors to 2050 is a separate assumption that this experiment doesn't make.
\item
  \textbf{Gas steam/turbine has no KEPCO cohort.} It takes the combined-cycle
  factor, so this experiment can't separate combined-cycle from steam/turbine
  plants. The evidence for that lumping is GCAM's own identical factors (Step
  6a). Gas steam/turbine is 3.9 TWh of the reference scenario's 2050
  generation.
\item
  \textbf{Why
  there is
  no upper
  bound
  for
  steel.}
  Filling
  hydrogen-based
  DRI with
  GCAM's
  own
  \texttt{EAF\ with\ DRI}
  factor
  instead
  makes
  net-zero
  steel
  NOx
  \emph{rise}
  40\%.
  That
  factor
  is
  itself
  the
  anomalous
  stock
  coefficient
  (11.66
  kg
  NOx/t;
  audit
  §B6.2),
  so we
  don't
  use it
  as a
  bound.
\item
  \textbf{Some
  weaknesses
  aren't
  tested.}
  Placeholder-driven
  sources
  (solvents,
  industrial
  processes)
  and
  refining
  are not
  part of
  the
  experiment,
  because
  they
  have no
  activity
  driver
  to apply
  a factor
  to.
\end{itemize}

\hypertarget{what-this-means-for-using-gcam-kaist}{%
\subsection*{5. What this means for using GCAM-KAIST \plainmark{}}\addcontentsline{toc}{subsection}{5. What this means for using GCAM-KAIST}\label{what-this-means-for-using-gcam-kaist}}

\textbf{Use
GCAM-KAIST's
activity;
don't use
its native
non-CO2
emissions.}
GCAM-KAIST
remains
the
project's
source of
projected
activity.
What it
can't
supply is
emission
factors
that
distinguish
technologies,
respond to
controls,
or cover
every
source.

\hypertarget{limitations-of-gcams-native-non-co2-approach-as-inherited-by-gcam-kaist}{%
\subsubsection*{5.1 Limitations of GCAM's native non-CO2 approach, as inherited by GCAM-KAIST}\addcontentsline{toc}{subsubsection}{5.1 Limitations of GCAM's native non-CO2 approach, as inherited by GCAM-KAIST}\label{limitations-of-gcams-native-non-co2-approach-as-inherited-by-gcam-kaist}}

\begin{longtable}[]{@{}
  >{\raggedright\arraybackslash}p{(\columnwidth - 2\tabcolsep) * \real{0.6389}}
  >{\raggedright\arraybackslash}p{(\columnwidth - 2\tabcolsep) * \real{0.3611}}@{}}
\toprule\noalign{}
\begin{minipage}[b]{\linewidth}\raggedright
Limitation
\end{minipage}
&
\begin{minipage}[b]{\linewidth}\raggedright
Evidence
\end{minipage} \\
\midrule\noalign{}
\endhead
\bottomrule\noalign{}
\endlastfoot
Base-year
emissions
are
calibrated
inventory
totals, so
agreement
with CAPSS
is
circular &
Step 4;
audit
§B5 \\
One factor
shared by
distinct
technologies
within an
inventory
total &
Step 6a;
audit
§B6.2 \\
Factors
swing
between
base years
& Step 6b;
audit
§B6.3 \\
Large
shares of
VOCs and
SOx tied
to a
placeholder,
not to
activity &
Step 6c;
audit
§B6.4 \\
Future
factors
frozen or
on a curve
shared by
both
scenarios;
no
pollution-control
lever &
Step 7;
audit
§B6.5 \\
Untagged
sources
and
technologies
count as
zero
(cement
process,
refining,
hybrids,
hydrogen
steel,
CCS) &
Step 3;
audit
§B4 \\
No primary
PM2.5 or
PM10 &
audit
§B2 \\
\end{longtable}

\hypertarget{the-power-sector-what-is-and-isnt-lumped}{%
\subsubsection*{5.2 The power sector: what is and isn't lumped}\addcontentsline{toc}{subsubsection}{5.2 The power sector: what is and isn't lumped}\label{the-power-sector-what-is-and-isnt-lumped}}

Power is
the
best-resolved
part of
GCAM-KAIST,
so it
needs a
precise
statement:

\begin{itemize}
\tightlist
\item
  \textbf{Each
  fuel has
  its own
  factor.}
  Coal,
  gas, oil
  and
  biomass
  power
  are
  separate,
  so a
  coal-to-gas
  switch
  does
  change
  GCAM's
  emissions
  in the
  right
  direction.
\item
  \textbf{Lumping
  happens
  \emph{within}
  a fuel.}
  Combined-cycle
  and
  steam/turbine
  gas
  plants
  share
  one
  factor
  per unit
  of fuel
  (Figure
  3). The
  coal
  IGCC and
  coal CCS
  variants
  have no
  activity
  in these
  runs,
  and gas
  with CCS
  runs
  only at
  small
  scale in
  2050
  (\texttt{nz}).
\item
  \textbf{How the base-year factors compare with KEPCO's measured factors.}
  We compare like years: GCAM-KAIST 2015 with KEPCO 2015, and GCAM-KAIST 2021
  with KEPCO 2021. The KEPCO factors are the hand-calculated,
  generation-weighted factors described in Section 4.2 and are not
  production-ready
  (\texttt{figures/\allowbreak{}data/\allowbreak{}table\_\allowbreak{}power\_\allowbreak{}factors.\allowbreak{}csv}).
\end{itemize}

\begin{longtable}[]{@{}p{0.22\linewidth}p{0.11\linewidth}p{0.10\linewidth}p{0.11\linewidth}p{0.10\linewidth}p{0.22\linewidth}@{}}
\toprule\noalign{}
Technology, pollutant & GCAM-KAIST 2015 & KEPCO 2015 & GCAM-KAIST 2021 & KEPCO 2021 & GCAM-KAIST ÷ KEPCO \\
\midrule\noalign{}
\endhead
\bottomrule\noalign{}
\endlastfoot
Coal, NOx & 0.458 & 0.534 & 0.150 & 0.126 & 0.86 (2015), 1.19 (2021) \\
Coal, SOx & 0.267 & 0.335 & 0.071 & 0.139 & 0.80 (2015), 0.51 (2021) \\
Gas combined cycle, NOx & 0.134 & 0.203 & 0.040 & 0.164 & 0.66 (2015), 0.24 (2021) \\
Gas steam/turbine, NOx & 0.201 & --- & 0.060 & --- & no KEPCO cohort \\
Oil, NOx & 0.509 & 0.571 & 0.320 & 0.301 & 0.89 (2015), 1.06 (2021) \\
Oil, SOx & 0.520 & 0.334 & 0.374 & 0.011 & 1.56 (2015), 34 (2021) \\
\end{longtable}

\emph{All values in kg/MWh. KEPCO cohorts: coal and gas combined cycle from
5--9 plants; oil from one plant.}

In short: \textbf{coal NOx and oil NOx track KEPCO within about 20\% in both
years. Coal SOx tracks in 2015 but is half of KEPCO's in 2021. Gas
combined-cycle NOx is too low in both years, 4.1× too low by 2021. Oil SOx is
far too high in 2021.} Even where a factor tracks, it is a fleet average. It
carries no information about which plants have which controls, and after
2021 it follows the shared curve rather than policy.

\hypertarget{beyond-this-project}{%
\subsubsection*{5.3 Beyond this project}\addcontentsline{toc}{subsubsection}{5.3 Beyond this project}\label{beyond-this-project}}

Because
these
properties
come from
GCAM's
standard
non-CO2
approach,
they apply
to any
study that
uses
GCAM's
native
non-CO2
emissions
to
evaluate
technology
transitions,
not only
to
GCAM-KAIST.
Whether
the same
lumping
holds in
GCAM-USA,
which has
its own US
emission
inputs,
has not
been
checked
here and
is the
subject of
a separate
note.

\hypertarget{our-approach-bottom-up-emission-factors}{%
\subsection*{6. Our approach: bottom-up emission factors \plainmark{}}\addcontentsline{toc}{subsection}{6. Our approach: bottom-up emission factors}\label{our-approach-bottom-up-emission-factors}}

The
project
builds a
non-power
(and
interim
power)
emission-factor
inventory
that pairs
\textbf{GCAM-KAIST
activity}
with
\textbf{technology-specific
Korean
emission
factors}.
It keeps
three legs
separate:

\begin{itemize}
\tightlist
\item
  \textbf{Activity:}
  which
  GCAM
  path and
  unit
  gives
  the
  physical
  quantity,
  for
  example
  tonnes
  of
  clinker
  or
  vehicle-km.
\item
  \textbf{Factor:}
  which
  CAPSS or
  Korean
  source
  supplies
  the
  emission
  factor
  for that
  quantity.
\item
  \textbf{Spatial:}
  where
  the
  emissions
  are
  placed
  on the
  grid.
\end{itemize}

How it
addresses
each
limitation:

\begin{longtable}[]{@{}
  >{\raggedright\arraybackslash}p{(\columnwidth - 2\tabcolsep) * \real{0.4651}}
  >{\raggedright\arraybackslash}p{(\columnwidth - 2\tabcolsep) * \real{0.5349}}@{}}
\toprule\noalign{}
\begin{minipage}[b]{\linewidth}\raggedright
GCAM
limitation
\end{minipage}
&
\begin{minipage}[b]{\linewidth}\raggedright
Bottom-up
remedy
\end{minipage} \\
\midrule\noalign{}
\endhead
\bottomrule\noalign{}
\endlastfoot
Circular
agreement
with CAPSS
& Factors
come from
source-level
evidence,
not from
dividing
inventory
totals.
Agreement
with CAPSS
in the
base years
then
becomes a
genuine
validation
test. \\
Shared
factors
across
technologies
& One
factor per
technology
and
control
configuration
(for
example
combined-cycle
vs
steam/turbine;
plants
with and
without
flue-gas
desulfurization). \\
Untagged
sources
count as
zero &
Every
activity
gets a
factor or
an
explicit,
recorded
gap; a
missing
factor
never
becomes
zero. \\
Placeholder
drivers &
Process
emissions
are tied
to
physical
activity:
clinker,
crude
throughput,
solvent
use. \\
No
pollution-control
lever &
Factors
can differ
by
scenario
where a
policy
changes
controls
or
standards. \\
No PM2.5
or PM10 &
CAPSS
factors
include
PM2.5,
PM10 and
TSP. \\
\end{longtable}

This
approach
has its
own costs:

\begin{itemize}
\tightlist
\item
  GCAM's
  units
  (EJ, Mt,
  passenger-km)
  must be
  converted
  to the
  factors'
  units
  (tonnes
  of
  clinker,
  vehicle-km,
  animal-years).
\item
  Korean
  factors
  must be
  collected
  and
  reviewed.
\item
  A
  spatial
  allocation
  must be
  built.
\end{itemize}

The
inventory
is a work
in
progress.
Current
status is
tracked in
{\texttt{roadmap.\allowbreak{}md}},
and the
method
will be
described
fully once
it is
complete.

\hypertarget{caveats-and-open-items}{%
\subsection*{7. Caveats and open items}\addcontentsline{toc}{subsection}{7. Caveats and open items}\label{caveats-and-open-items}}

\begin{itemize}
\tightlist
\item
  \textbf{KAIST
  inputs.}
  GCAM-KAIST's
  input
  files
  were not
  examined.
  The
  provenance
  argument
  (Step 4)
  rests on
  its
  outputs
  matching
  stock
  GCAM v9.
  Requested
  from
  KAIST:
  the run
  configurations
  and
  emissions
  input
  files.
\item
  \textbf{CEDS release.} CEDS's code confirms the scaling to CAPSS. But the
  release GCAM v9 uses isn't recorded (inferred to be 2024 or later), and CEDS's
  second Korea scaling step (EDGAR-HTAP, 2000--2018) covers 2015 but not 2021
  (audit §B5.3).
\item
  \textbf{Refinery
  emissions.}
  That
  they sit
  in
  \texttt{other\ industrial\ energy\ use\ /\allowbreak{}\ refined\ liquids}
  is an
  inference
  (Step 5;
  audit
  §B5.5).
\item
  \textbf{Mapping
  search.}
  Choosing
  the best
  of
  49,152
  mappings
  flatters
  the fit;
  only the
  ``remapping
  can't
  create
  agreement''
  conclusion
  is drawn
  from it.
\item
  \textbf{Policy
  experiment.}
  Illustrative,
  with
  factors
  that are
  not
  production-ready
  (§4.4).
\item
  \textbf{Region.}
  Only
  South
  Korea
  was
  examined.
\end{itemize}

\hypertarget{figures-tables-and-reproduction}{%
\subsection*{8. Figures, tables and reproduction}\addcontentsline{toc}{subsection}{8. Figures, tables and reproduction}\label{figures-tables-and-reproduction}}

All
figures
and tables
are
produced
by one
script,
{\texttt{figures/\allowbreak{}make\_\allowbreak{}figures.\allowbreak{}py}},
run from
the repo
root. Each
figure's
plotted
numbers
are saved
in
\texttt{figures/\allowbreak{}data/\allowbreak{}}
(the table
view). The
script
reads
local
diagnostics
under
\texttt{results/\allowbreak{}}
and
\texttt{data/\allowbreak{}},
which are
untracked,
so it can
only be
re-run
where
those
files
exist.

\begin{longtable}[]{@{}
  >{\raggedright\arraybackslash}p{(\columnwidth - 4\tabcolsep) * \real{0.2756}}
  >{\raggedright\arraybackslash}p{(\columnwidth - 4\tabcolsep) * \real{0.3622}}
  >{\raggedright\arraybackslash}p{(\columnwidth - 4\tabcolsep) * \real{0.3622}}@{}}
\toprule\noalign{}
\begin{minipage}[b]{\linewidth}\raggedright
Item
\end{minipage}
&
\begin{minipage}[b]{\linewidth}\raggedright
File
\end{minipage}
&
\begin{minipage}[b]{\linewidth}\raggedright
Data
\end{minipage} \\
\midrule\noalign{}
\endhead
\bottomrule\noalign{}
\endlastfoot
Figure 1 &
\texttt{figures/\allowbreak{}fig01\_\allowbreak{}capss\_\allowbreak{}vs\_\allowbreak{}gcam\_\allowbreak{}2021.\allowbreak{}png}
(2015:
\texttt{.\allowbreak{}.\allowbreak{}.\allowbreak{}\_\allowbreak{}2015.\allowbreak{}png})
&
\texttt{data/\allowbreak{}fig01\_\allowbreak{}capss\_\allowbreak{}vs\_\allowbreak{}gcam\_\allowbreak{}\{2015,2021\}.\allowbreak{}csv} \\
Figure 2 &
\texttt{figures/\allowbreak{}fig02\_\allowbreak{}capss\_\allowbreak{}agreement.\allowbreak{}png}
&
\texttt{data/\allowbreak{}fig02\_\allowbreak{}capss\_\allowbreak{}agreement.\allowbreak{}csv} \\
Figure 3 &
\texttt{figures/\allowbreak{}fig03\_\allowbreak{}gas\_\allowbreak{}cc\_\allowbreak{}vs\_\allowbreak{}steamct.\allowbreak{}png}
&
\texttt{data/\allowbreak{}fig03\_\allowbreak{}gas\_\allowbreak{}cc\_\allowbreak{}vs\_\allowbreak{}steamct.\allowbreak{}csv} \\
Figure 4 &
\texttt{figures/\allowbreak{}fig04\_\allowbreak{}mobile\_\allowbreak{}vs\_\allowbreak{}stationary.\allowbreak{}png}
&
\texttt{data/\allowbreak{}fig04\_\allowbreak{}mobile\_\allowbreak{}vs\_\allowbreak{}stationary.\allowbreak{}csv} \\
Figure 5 &
\texttt{figures/\allowbreak{}fig05\_\allowbreak{}base\_\allowbreak{}year\_\allowbreak{}swing.\allowbreak{}png}
&
\texttt{data/\allowbreak{}fig05\_\allowbreak{}base\_\allowbreak{}year\_\allowbreak{}swing.\allowbreak{}csv} \\
Figure 6 &
\texttt{figures/\allowbreak{}fig06\_\allowbreak{}placeholder\_\allowbreak{}shares.\allowbreak{}png}
&
\texttt{data/\allowbreak{}fig06\_\allowbreak{}placeholder\_\allowbreak{}shares.\allowbreak{}csv} \\
Figure 7 &
\texttt{figures/\allowbreak{}fig07\_\allowbreak{}coefficient\_\allowbreak{}trajectories.\allowbreak{}png}
&
\texttt{data/\allowbreak{}fig07\_\allowbreak{}coefficient\_\allowbreak{}trajectories.\allowbreak{}csv} \\
Figure 8 &
\texttt{figures/\allowbreak{}fig08\_\allowbreak{}policy\_\allowbreak{}power.\allowbreak{}png}
&
\texttt{data/\allowbreak{}fig08\_\allowbreak{}policy\_\allowbreak{}power.\allowbreak{}csv} \\
Figure 9 &
\texttt{figures/\allowbreak{}fig09\_\allowbreak{}policy\_\allowbreak{}steel.\allowbreak{}png}
&
\texttt{data/\allowbreak{}fig09\_\allowbreak{}policy\_\allowbreak{}steel.\allowbreak{}csv} \\
Figure 10
&
\texttt{figures/\allowbreak{}fig10\_\allowbreak{}policy\_\allowbreak{}road.\allowbreak{}png}
&
\texttt{data/\allowbreak{}fig10\_\allowbreak{}policy\_\allowbreak{}road.\allowbreak{}csv} \\
Mapping-search
table
(Step 5) &
&
\texttt{data/\allowbreak{}table\_\allowbreak{}mapping\_\allowbreak{}search.\allowbreak{}csv} \\
Power-factor
table
(§5.2) & &
\texttt{data/\allowbreak{}table\_\allowbreak{}power\_\allowbreak{}factors.\allowbreak{}csv} \\
\end{longtable}

Upstream
evidence
scripts
are listed
in audit
§B9. The
experiment's
script and
outputs
are in
\texttt{results/\allowbreak{}diagnostics/\allowbreak{}policy\_\allowbreak{}experiment/\allowbreak{}}.
